% 调和函数（综述）
% license CCBYSA3
% type Wiki

本文根据 CC-BY-SA 协议转载翻译自维基百科\href{https://en.wikipedia.org/wiki/Harmonic_function}{相关文章}。

在数学、数理物理以及随机过程理论中，**调和函数（harmonic function）**是指满足拉普拉斯方程（Laplace’s equation）的函数。具体来说，若

[
f : U \to \mathbb{R}
]

是定义在欧几里得空间 (\mathbb{R}^n) 的一个**二阶连续可微函数**（即 (f \in C^2(U))），其中 (U) 是 (\mathbb{R}^n) 的一个开子集，那么若该函数在 (U) 内处处满足

[
\frac{\partial^2 f}{\partial x_1^2} + \frac{\partial^2 f}{\partial x_2^2} + \cdots + \frac{\partial^2 f}{\partial x_n^2} = 0,
]

则称 (f) 为一个调和函数。
这一方程通常可写作

[
\nabla^2 f = 0 \quad \text{或} \quad \Delta f = 0.
]

---

### “Harmonic（调和）”一词的词源

“调和函数（harmonic function）”中的形容词“harmonic”（调和的、谐和的）源自**“谐振运动”**（harmonic motion）——例如一根被拉紧的弦上某一点的振动。描述这种运动的微分方程的解可用**正弦函数**和**余弦函数**表示，因此这些函数被称为**谐波函数（harmonics）**。

**傅里叶分析（Fourier analysis）**即是在单位圆上将任意函数展开为这些谐波函数的级数。将这一思想推广到更高维空间（即单位 (n) 维球面上的类似结构），便得到所谓的**球谐函数（spherical harmonics）**。这些函数恰好满足拉普拉斯方程。随着时间推移，“harmonic”一词逐渐被推广，用来指代**所有满足拉普拉斯方程的函数**【1】。
